% !TeX root = main.tex
\section{Twisted sheaves}
\label{sec:quasi-coher-sheav}
In this section, $\Delta$ will denote a diagonalizable group scheme (over $\Spec \mathbb{Z}$)
with character group~$A := \Hom_{\Spec\mathbb{Z}}(\Delta, \GGm)$.
Recall that $A$ is a finitely generated abelian group.
We will usually use multiplicative notation for $A$,
except when we identify it with something like the additive group of $\mathbb{Z}$.

Let $S$ be an algebraic stack.
By a $\Delta$-gerbe on $S$ we mean a $\Delta_S$-gerbe on
$S$.
Fix a $\Delta$-gerbe $\alpha\colon X \to S$.
Note that $X$ is isomorphic to $\BB \Delta_S$ locally on $S$.
In particular, $X$ is an algebraic stack,
and it is reasonable to talk about quasi-coherent sheaves on~$X$.

\begin{remark}
\label{rem-gerbe-ff}
Let $\alpha\colon X \to S$ be a $\Delta$-gerbe.
Note that the structure morphism $\alpha$ is smooth and surjective.
Indeed, this is true for any algebraic gerbe
\cite[Proposition~A.2]{bergh2017}.
Moreover, assume that $Y$ is a $\Delta'$-gerbe over $S$
and let $\rho\colon X \to Y$ be a morphism over $S$ with induced morphism $\varphi\colon \Delta_S \to \Delta'_S$ on bands.
Then $\rho$ factors as $X \to X\times\BB \Delta' \to \varphi_*X \cong Y$ by Construction~\ref{cons-gerbe-pf}
combined with Proposition~\ref{gerbe-functorial-group}.
Note that the first morphism in the factorization is a
$\Delta'$-torsor 
and that the second morphism is a gerbe banded by~$\Delta$.
In particular, the morphism $\rho$ is faithfully flat.
See also Remark~\ref{rem-gerbe-concentrated} for cohomological properties of $\rho$.
\end{remark}

\begin{definition}
\label{def-homogeneous-sheaf}
Let $\mathcal{F}$ be an object in $\Mod(X_\fppf, \mathcal{O}_X)$,
and let $\chi\colon \Delta_S \to \GGms{S}$ be a character.
Note that $\Delta_X$ acts on $\mathcal{F}$ from the right via the inertial action \eqref{eq-inertial-action}
and the banding (see Definition~\ref{def-banding}) of $X$,
and that $\GGm$ acts on $\mathcal{F}$ from the left via the inclusion $\GGms{X} \subset \mathcal{O}_X$.
We define the \emph{$\chi$-homogeneous subsheaf} $\mathcal{F}_\chi$ of $\mathcal{F}$ by
\begin{equation}
  \label{eq:chi-subsheaf}
  \mathcal{F}_\chi = \{s \in \mathcal{F} \mid 
  \text{$s\cdot\gamma = \chi(\gamma)s$ for all $\gamma
  \in \Delta_X$}\}.
\end{equation}
An object $\mathcal{F}$ of $\Mod(X_\fppf, \mathcal{O}_X)$ is
called \emph{$\chi$-homogeneous}
or \emph{homogeneous of degree $\chi$}
provided that $\mathcal{F} = \mathcal{F}_\chi$,
and we denote the full subcategory of $\chi$-homogeneous objects in $\Qcoh(X_\fppf, \mathcal{O}_X)$
by $\Qcoh_\chi(X_\fppf, \mathcal{O}_X)$.
\end{definition}

\begin{remark}
\label{rem:def-Qcoh-chi}
Note that the expression
\eqref{eq:chi-subsheaf} does not make sense in the lisse--\'etale
topos unless $\Delta_X$ is smooth over $X$.
When we talk about homogeneous sheaves in $\Qcoh(X_\liset, \mathcal{O}_X)$,
we implicitly transport the subcategory
$\Qcoh_\chi(X_\fppf, \mathcal{O}_X) \subset \Qcoh(X_\fppf, \mathcal{O}_X)$
to a subcategory $\Qcoh_\chi(X_\liset, \mathcal{O}_X) \subset \Qcoh(X_\liset, \mathcal{O}_X)$
under the equivalence obtained by restriction.
We will usually simply write $\Qcoh_\chi(X)$ for any of these subcategories. 
\end{remark}

\begin{definition}[Twisted sheaves]
\label{def:twisted}
Let $\alpha\colon X \to S$ be a $\GGm$-gerbe and let $\id\colon \GGm \to \GGm$
be the identity character.
An \emph{$\alpha$-twisted quasi-coherent sheaf on $S$} is simply an object of $\Qcoh_\id(X)$.
We will often use the notation $\Qcoh_{\alpha}(S)$ instead
of $\Qcoh_\id(X)$.
\end{definition}

\begin{remark}
\label{rem-twisted-abuse}
Since the category $\Qcoh_{\alpha}(S)$, up to equivalence,
obviously depends only on the class of $\alpha$ in $\HH^2(S_\fppf, \GGm)$,
we will sometimes abuse notation and write $\Qcoh_{\alpha}(S)$
even if $\alpha$ is just a cohomology class in $\HH^2(S_\fppf, \GGm)$.
\end{remark}

\begin{remark}
\label{rem-twisted-2-cocycle}
Twisted sheaves may also be described using 2-cocycles.
This approach was taken by C{\u{a}}ld{\u{a}}raru in \cite[Definition~1.2.1]{caldararu2000}.
That his definition is essentially equivalent to the definition above was proved by
Lieblich in~\cite[Proposition~2.1.3.11]{lieblich2004}. 
\end{remark}

We now give a precise formulation of the well-known fact that
any quasi-coherent sheaf on a suitable gerbe
decomposes into its homogeneous subsheaves (c.f.~\cite[Proposition~2.2.1.6]{lieblich2004}).

\begin{theorem}
  \label{t-qc-gerbe-split}
  Let $S$ be an algebraic stack and let $\alpha\colon X \to S$ be a
  gerbe 
  banded by a diagonalizable group $\Delta$.
  Then $\Qcoh_\chi(X)$ is
  a Serre subcategory (see~\cite[\sptag{02MO}]{stacks-project}) of $\Qcoh(X)$
  for any character $\chi$ of $\Delta_S$,
  and taking the coproduct gives an equivalence
  \begin{equation}
    \label{eq-eq-qc-gerbe}
    \prod_{\chi \in A}\Qcoh_\chi(X) \sira \Qcoh(X),
    \qquad
    (\mathcal{F}_\chi)_\chi \mapsto \bigoplus_{\chi \in A}
    \mathcal{F}_\chi, 
  \end{equation}
  of abelian categories, where $A$ denotes the character group of $\Delta$.
  
  Furthermore, assume that $\mathcal{F}_\chi$ and $\mathcal{G}_\psi$ are quasi-coherent
  $\mathcal{O}_X$-modules which are homogeneous for $\Delta_S$-characters $\chi$ and $\psi$, respectively.
  Then the quasi-coherent $\mathcal{O}_X$-modules
  \begin{equation}
    \label{eq-qc-monoidal}
    \mathcal{F}_\chi \otimes \mathcal{G}_\psi, \qquad
    \sheafHom(\mathcal{G}_\psi, \mathcal{F}_\chi)
  \end{equation}
  are $\chi\psi$- and $\chi\psi^{-1}$-homogeneous,
  respectively. 
\end{theorem}

\begin{proof}
Throughout the proof,
we work with sheaves in $X_\fppf$,
which we view as the topos of sheaves on the site of affine schemes over $X$.
Fix a quasi-coherent sheaf $\mathcal{F}$ of $\mathcal{O}_X$-modules.

  Let $x$ be an object of $X$ lying over an affine scheme $\Spec A$.
  Then $\mathcal{F}(x)$ is an $A$-module with an $A$-linear right action of the group 
  $\Delta(A):=\Delta(\Spec A)$ coming from the inertial action via the banding.
  Given a morphism $f \colon y \ra x$ in $X$ lying over a
  morphism $\Spec \phi \colon \Spec B \ra \Spec A$,
  the corresponding morphism $\mathcal{F}(f) \colon \mathcal{F}(x) \ra \mathcal{F}(y)$ 
  is equivariant with respect to $\Delta(A) \ra \Delta(B)$.
  Moreover, since $\mathcal{F}$ is quasi-coherent,
  it induces an isomorphism $B \otimes_{\phi, A} 
  \mathcal{F}(x) \sira \mathcal{F}(y)$ of $B$-modules.

  By restricting $\mathcal{F}$ to $\Spec A$ along $x\colon \Spec A \to X$ (see \cite[\sptag{06W9}]{stacks-project}),
  we see that the $A$-module $\mathcal{F}(x)$ can be viewed as
  a right representation of the diagonalizable $A$-group scheme
  $\Delta_{\Spec A}$ in the sense of \cite[I.2.7]{Jantzen}.
  Hence it decomposes as 
  $\mathcal{F}(x)=\bigoplus_{\chi \in A}\mathcal{F}(x)_\chi$ by 
  \cite[I.2.11]{Jantzen}.

  If the object $x$ varies, all these decompositions 
  are compatible and combine into the decomposition 
  $\mathcal{F}=\bigoplus_{\chi \in A} \mathcal{F}_\chi$ where
  $\mathcal{F}_\chi$ is the 
  $\chi$-homogeneous subsheaf of $\mathcal{F}$. 
  This shows that the functor \eqref{eq-eq-qc-gerbe} is
  essentially surjective.
  Similarly, we get a decomposition of morphisms between pairs of quasi-coherent modules
  (cf. \cite[I.2.11]{Jantzen}),
  which shows that \eqref{eq-eq-qc-gerbe} is fully faithful.
  Moreover, any $\Delta_S$-character $\chi$ clearly gives an exact functor $\mathcal{F} \mapsto \mathcal{F}_\chi$.
  In particular, $\Qcoh_\chi(X)$ is a Serre subcategory of $\Qcoh(X)$.

  The statement that the tensor product of a $\chi$-homogeneous
  quasi-coherent module with a $\psi$-homogeneous quasi-coherent
  module is $(\chi+\psi)$-homogeneous is obvious, and
  the corresponding claim for the internal Hom-functor is then a
  formal consequence.
\end{proof}

\begin{remark}
\label{rem-coprod-prod}
Note that the canonical morphism
$
\bigoplus_{\chi\in A} \mathcal{F}_\chi
\to
\prod_{\chi\in A} \mathcal{F}_\chi
$
is an isomorphism even if the character group
$A$ is not finite. Here the product is taken in $\Qcoh(X)$.
This is a simple consequence of 
the equivalence 
\eqref{eq-eq-qc-gerbe} in
Theorem~\ref{t-qc-gerbe-split}.
\end{remark}

\begin{proposition}
\label{p:pull-back-grading}
Let $S$ be an algebraic stack and 
$\alpha\colon X \to S$ a gerbe banded by a diagonalizable group
$\Delta$.
Let $f\colon T \to S$ be a morphism of algebraic stacks,
and let $g\colon Y \to X$ denote the base change of $f$ along
$\alpha$.
Then the functors
$$
g^*\colon \Qcoh(X) \to \Qcoh(Y), \qquad
g_*\colon \Qcoh(Y) \to \Qcoh(X)
$$
respect the decompositions of $\Qcoh(X)$ and~$\Qcoh(Y)$
into homogeneous objects given in Theorem~\ref{t-qc-gerbe-split}.
\end{proposition}
\begin{proof}
As a formal consequence of adjunction,
it is enough to verify that one of the functors $g^*$ and~$g_*$ preserves the decomposition.
The functor $g^*$ is obtained by restriction of the functor
$g^*\colon\Mod(S_\fppf, \mathcal{O}_S) \to \Mod(T_\fppf, \mathcal{O}_T)$,
which in turn is just given by restriction of the topos $S_\fppf$.
In this setting the statement is obvious from the definition of a
homogeneous sheaf
(see Definition~\ref{def-homogeneous-sheaf}).
\end{proof}

\begin{proposition}
  \label{t:equivalence-epi}
  Let $X$ and $X'$ be gerbes over $S$ banded by diagonalizable groups $\Delta$ and~$\Delta'$,
  respectively, and let $\rho\colon X \to X'$ be a morphism over~$S$ with induced morphism $\varphi\colon \Delta_S \to \Delta'_S$ on the bands.
  Fix a character $\chi'\colon \Delta'_S \to \GGms{S}$, and let $\chi =
  \chi' \circ \varphi$.
  Then the pull-back functor $\rho^*$ takes $\chi'$-homogeneous
  objects to $\chi$-homogeneous 
  objects and induces an equivalence
  \begin{equation}
    \label{eq-rho-restricted}
    \rho^*\colon 
    \Qcoh_{\chi'}(X') \sira \Qcoh_{\chi}(X)
  \end{equation}
  of categories.
  In particular, if $X' = S$ then $\rho^*\colon \Qcoh(S) \to \Qcoh(X)$ is fully faithful with 
  essential image consisting of those objects that are
  homogeneous for the trivial character. 
\end{proposition}

\begin{proof}
By a similar argument as in the proof of Proposition~\ref{p:pull-back-grading},
we see that $\rho^*$
takes $\chi'$-homogeneous objects to 
$\chi$-homogeneous objects.
First we note that the special case in the last
sentence in the statement of the proposition follows
directly from 
Proposition~\ref{p-gerbe-pb-ff}.

In the general case, we consider the fiber product $X\times_S X'$
and its projections $p$ and~$q$ to $X$ and~$X'$.
Note that $X\times_S X'$ is a $\Delta\times\Delta'$-gerbe over~$S$
and that $p$ and~$q$ induce the projections to $\Delta$ respectively~$\Delta'$
on the bands.

The projection $p$ is a gerbe banded by the kernel of the
projection $\Delta\times\Delta' \to \Delta$,
which is isomorphic to~$\Delta'$.
In particular, $p^*$ is fully faithful with essential image consisting of sheaves
on which this kernel acts trivially.
Similar statements hold for~$q$.
Hence it follows from the decomposition in Theorem~\ref{t-qc-gerbe-split}
that we get equivalences
\begin{align}
  \label{eq:equiv-p}
  p^* \colon \Qcoh_\chi(X)
  & \rightleftarrows \Qcoh_{(\chi,e')}(X \times_S X'/S)\colon p_*,\\
  \label{eq:equiv-q}
  q^* \colon \Qcoh_{\chi'}(X')
  & \rightleftarrows \Qcoh_{(e, \chi')}(X \times_S X'/S)\colon q_*,
\end{align}
where $e$ and~$e'$ denote the trivial characters of $\Delta$ and~$\Delta'$,
respectively.

Note that the gerbe $p$ has a trivialization~$\tau = (\id, \rho)$.
In particular, there exists a line bundle $\mathcal{L}$ on $X\times_S X'$ such
that
\begin{enumerate}
\item
\label{it-cond-1}
$\mathcal{L}$ is $\chi'$-homogeneous with respect to the decomposition induced by~$p$;
\item
\label{it-cond-2}
$\tau^*\mathcal{L} \cong \mathcal{O}_X$.
\end{enumerate}
Since $\mathcal{L}$ is a line bundle,
it must also be homogeneous for some character $(\psi, \psi')$ of $(\Delta\times\Delta')_S$
with respect to the gerbe $X\times_S X' \to S$.
By \ref{it-cond-1}, we have $\psi' = \chi'$.
Furthermore, since $\tau$ induces the morphism $(\id, \varphi)$ on bands,
we have $\tau^*\mathcal{L} \in \Qcoh_{\psi\chi}(X)$.
This forces $\psi = \chi^{-1}$ by \ref{it-cond-2},
so $\mathcal{L} \in \Qcoh_{(\chi^{-1}, \chi')}(X\times_S X'/S)$.
In particular, it follows from the equivalences \eqref{eq:equiv-p} and~\eqref{eq:equiv-q} and the statement about
the monoidal structure in Theorem~\ref{t-qc-gerbe-split} that we
get an equivalence
$$
p_*(\mathcal{L}^\vee\otimes q^*(-))\colon \Qcoh_{\chi'}(X') \sira \Qcoh_{\chi}(X).
$$
Furthermore, the restriction of the natural transformation 
\begin{multline*}
p_*(\mathcal{L}^\vee\otimes q^*(-)) \cong \tau^*p^*p_*(\mathcal{L}^\vee\otimes q^*(-)) \to\\
\tau^*(\mathcal{L}^\vee\otimes q^*(-)) \cong \tau^*(\mathcal{L}^\vee)\otimes \tau^*q^*(-) \cong \rho^*
\end{multline*}
to $\Qcoh_{\chi'}(X')$ is an isomorphism,
which concludes the proof.
\end{proof}

Now the following corollaries are easy consequences of Theorem~\ref{t-qc-gerbe-split},
Proposition~\ref{p:pull-back-grading} and~Proposition~\ref{t:equivalence-epi}.
\begin{corollary}
\label{cor:twisted-power}
Let $S$ be an algebraic stacks and let $\alpha\colon X \to S$ be a gerbe banded by $\Delta$,
and choose a character $\chi\colon \Delta_S \to \GGms{S}$.
Then we have an equivalence
\begin{equation}
\label{eq-pb-character}
\Qcoh_{\chi_*\alpha}(S)\cong \Qcoh_\chi(X) 
\end{equation}
of categories, where $\chi_*\alpha$ is defined as in Construction~\ref{cons-gerbe-pf}.
In the particular case when $\Delta = \GGm$,
this gives the equivalence
\begin{equation}
\label{eq-pb-power}
\Qcoh_{\alpha^d}(S)\cong \Qcoh_{d}(X) 
\end{equation}
for any $d \in \mathbb{Z}$, where $\alpha^d$ is defined as in Construction~\ref{cons-gerbe-operations}.
\end{corollary}
\begin{proof}
The equivalence \eqref{eq-pb-character} is the special case of Proposition~\ref{t:equivalence-epi} where
we take $\rho$ to be $X \to \chi_* X$ as in Construction~\ref{cons-gerbe-pf}.
The equivalence \eqref{eq-pb-power} is the special case of \eqref{eq-pb-character},
where we take $\chi = p_n$ as in Construction~\ref{cons-gerbe-operations}.
\end{proof}

\begin{corollary}
\label{cor:twisted-tensor}
Let $S$ be an algebraic stack and let $\alpha\colon X \to S$ and~$\beta\colon Y \to S$ be $\GGm$-gerbes.
Then the tensor product and the internal Hom-functor on $\Qcoh(X\times_S Y)$ induce functors
\begin{align*}
\otimes & \colon \Qcoh_{\alpha}(S) \times \Qcoh_{\beta}(S) \to \Qcoh_{\alpha\beta}(S), \\
\sheafHom & \colon \Qcoh_{\alpha}(S)^\opp \times \Qcoh_{\beta}(S)
\to \Qcoh_{\alpha^{-1}\beta}(S).
\end{align*}
Furthermore, if $f\colon T \to S$ is an arbitrary morphism of algebraic stacks,
then we get an induced pair of adjoint functors
\begin{equation*}
  f^* \colon \Qcoh_{\alpha}(S) \to \Qcoh_{f^*\alpha}(T),
  \qquad
  f_* \colon \Qcoh_{f^*\alpha}(T) \to \Qcoh_{\alpha}(S).
\end{equation*}
Recall our conventions regarding $f_*$ and $\sheafHom$ from Notation~\ref{not-sheaf-op}.
\end{corollary}
\begin{proof}
The product $X\times_S Y$ has the structure of a $\GGm\times \GGm$-gerbe
such that the projections on the factors $X$ and $Y$ induce the projections on the factors
on the bands.
By applying Proposition~\ref{t:equivalence-epi} to the two projections $\GGm\times \GGm \to \GGm$,
we get equivalences
$$
\Qcoh_{\alpha}(S) \cong \Qcoh_{(1, 0)}(X\times_S Y), \qquad
\Qcoh_{\beta}(S) \cong \Qcoh_{(0, 1)}(X\times_S Y),
$$
where we have identified the character group of $\GGm$ with~$\mathbb{Z}$.
By applying Proposition~\ref{t:equivalence-epi} to the multiplication and inversion maps
given in Construction~\ref{cons-gerbe-operations},
we obtain the equivalences
$$
\Qcoh_{\alpha\beta}(S) \cong \Qcoh_{(1, 1)}(X\times_S Y), \qquad
\Qcoh_{\alpha^{-1}\beta}(S) \cong \Qcoh_{(-1, 1)}(X\times_S Y),
$$
respectively.
Now the functors $\otimes$ and $\sheafHom$ are obtained from the monoidal structure of $\Qcoh(X\times_S Y)$
described in Theorem~\ref{t-qc-gerbe-split}.
The statement about the functors $f^*$ and $f_*$ is an obvious consequence of Proposition~\ref{p:pull-back-grading}.
\end{proof}

%%% Local Variables: 
%%% mode: latex
%%% TeX-master: "main"
%%% End: 
